\section{Cancellation along a path}
\label{sec:spectrum}

The comparison of the two observation copies led to sums of the edge
differences \(D_e(t)\).  Cauchy--Schwarz bounds the second moment of a
sum over \(r\) edges by \(r^2S(t)\), where \(S(t)=\E D_e(t)^2\).
We now improve the factor \(r^2\) to \(r\) by proving that the
correlations of these differences are summable along a path.
The proof begins with local correlations under one observation law
and passes to the joining through the closed space generated by
differences of marginal functions.  Although only equality remains
in the comparison, this path estimate holds throughout
\(0<\theta\le b^{-1/2}\).

Here and below \(\nu\) is the observation marginal of the original
joint law \(\mathsf Q\), namely the law of the entire array
\[
 X=(X_v)_{v\in\mathcal V}
 \quad\hbox{in}\quad
 \mathsf X=C_0([0,\infty),\mathbb R)^{\mathcal V}.
\]
Here \(C_0\) again means continuous paths starting at zero, with
the topology of uniform convergence on compact time intervals.
The vertex and nonbacktracking operator geometry used in this argument
is related to the correlation arguments of Backhausz, Szegedy, and
Vir\'ag~\cite[Theorem~2.2 in arXiv:1305.6784v1]{BSV}
and Backhausz, Gerencs\'er, Harangi, and
Vizer~\cite[Theorems~4.1--4.2 in arXiv:1603.08423v2]{BGHV}.  The proof below supplies the precise transfer to
differences of marginal observation functions.

\begin{proposition}[Path variance]
\label{prop:path-variance}
Let \(d\ge3\), \(b=d-1\), and \(0<\theta\le b^{-1/2}\).
Let \(\pi\) be a probability law on \(\mathsf X^2\) with both
marginals equal to \(\nu\), invariant under the simultaneous action
of \(\Aut(T_d)\).  Suppose that \(g=(g_e)\) assigns a measurable
real function of \(\mathsf X\) to every directed edge, with
\(\E_\nu g_e^2<\infty\), and that for each fixed
\(\gamma\in\Aut(T_d)\) and directed edge \(e\),
\[
 g_{\gamma e}(\gamma x)=g_e(x)
 \quad\text{for \(\nu\)-almost every \(x\)}.
\]
Here \(\gamma(i,j)=(\gamma i,\gamma j)\).
For the two coordinate arrays \(X,X'\) under \(\pi\), set
\[
 D_e=g_e(X)-g_e(X'),\qquad S=\E_\pi D_e^2.
\]
The value of \(S\) is independent of the directed edge.  There is a
constant \(C_d<\infty\), depending only on \(d\), such that for every
integer \(r\ge1\) and every simple path \(v_0,\ldots,v_r\),
\begin{equation}
\label{eq:path-variance}
 \begin{aligned}
 \E_\pi\left[\left(\sum_{i=1}^rD_{(v_{i-1},v_i)}\right)^2\right]
     &\le C_d rS,\\
 \E_\pi\left[\left(\sum_{i=1}^rD_{(v_i,v_{i-1})}\right)^2\right]
     &\le C_d rS.
 \end{aligned}
\end{equation}
In particular, take \(\pi\) to be the \((X,X')\)-marginal of the
joining in Proposition~\ref{prop:compatible-copies}.  Then this holds for the
differences \(D_e(t)\) in \eqref{eq:log-discrepancy-def} at each fixed
finite time, with \(S=S(t)\).
\end{proposition}

We prove the proposition in three steps.  First, local observation
functions have exponentially small correlations across a long
connecting path.  This bounds the self-adjoint neighbor average on
centered vertex functions.  Second, the bound transfers to a closed
space of differences inside the joining.  Finally, head and tail lifts
of that vertex space control the nonbacktracking edge average.  Its
powers are summable, which gives the linear variance in
\eqref{eq:path-variance}.  All equivariance identities in this section
are understood separately for each fixed automorphism, as in the
proposition.

\subsection{The vertex bound for one observation law}

We shall use the same Hilbert-space notation for several invariant
configuration laws.  For such a law \(\eta\), a \emph{vertex
assignment} is a family \(f=(f_v)_{v\in\mathcal V}\) of measurable
real functions satisfying
\[
 f_{\gamma v}(\gamma x)=f_v(x)
 \quad\text{\(\eta\)-almost surely, for each fixed \(\gamma,v\)}.
\]
Let \(\mathcal H_{\mathrm v}(\eta)\) be the space of square-integrable
vertex assignments, modulo almost sure equality at each vertex, with
\[
 \langle f,k\rangle_{\mathrm v,\eta}=\E_\eta[f_ok_o].
\]
Transitivity and invariance make this independent of the vertex \(o\).
The analogous space \(\mathcal H_{\mathrm e}(\eta)\) consists of
assignments to directed edges, with inner product at any one directed
edge.  These spaces are Hilbert spaces: a Cauchy sequence has an
\(L^2\) limit at each index, and the equivariance identity for any
fixed automorphism passes to the limits because that automorphism
preserves \(\eta\).  There are countably many indices, so equality in
the indicated norm is equality of assignments.  We use real spaces;
the usual spectral theorem applies after complexification if needed.

On \(\mathcal H_{\mathrm v}(\eta)\) define the neighbor average
\[
 (\mathsf A_\eta f)_v=\frac1d\sum_{u\sim v}f_u.
\]
Jensen's inequality and invariance give \(\|\mathsf A_\eta\|\le1\).
For adjacent \(o,u\), an automorphism interchanging them gives
\(\E_\eta[f_ok_u]=\E_\eta[f_uk_o]\).  Averaging this identity over
the neighbors of \(o\) proves that \(\mathsf A_\eta\) is self-adjoint.
It preserves the closed subspace of centered assignments
\(\E_\eta f_o=0\).  This subspace is reducing: if \(y\) is
orthogonal to it, then
\(\langle \mathsf A_\eta y,f\rangle
=\langle y,\mathsf A_\eta f\rangle=0\)
for every centered \(f\).  Thus the restriction to that subspace
is again self-adjoint.

\begin{lemma}[Marginal vertex spectrum]
\label{lem:marginal-spectrum}
For the full observation law \(\nu\), the spectrum of \(\mathsf A_\nu\)
on centered vertex assignments is contained in \([-\rho,\rho]\), where
\begin{equation}
\label{eq:spectrum-rho}
 \rho=\frac{2\sqrt b}{d}<1.
\end{equation}
\end{lemma}

\begin{proof}
We first prove a correlation bound for an assignment \(f\) that is
centered and measurable, at vertex \(v\), from the entire observation
paths at the vertices in a fixed ball \(B_R(v)\).  Square
integrability is enough; write \(\|f\|^2=\E_\nu f_o^2\).
Write \(\E_{\mathrm{pl}}\) for expectation in the planted
spin-and-Brownian construction, whose observation marginal is \(\nu\).

Suppose that \(\dist(u,v)=k>2R\), and let \(a\) and \(c\) be the
respective endpoints of the path joining \(B_R(u)\) to \(B_R(v)\).
Then \(\dist(a,c)=k-2R\).  In the planted broadcast construction,
conditional on \((\sigma_a,\sigma_c)\), the observation arrays in
these two balls are independent.  The conditional law in the first
ball depends only on \(\sigma_a\), and the law in the second only on
\(\sigma_c\).  To see this, root the broadcast at \(a\).  Once the
endpoint spins are fixed, the edge multipliers inside the two balls
are taken from disjoint branches off the connecting path.  They are
independent of that path and of one another, and the Brownian paths
at the two disjoint sets of vertices are independent as well.
The assertion for products of square-integrable functions follows
from its bounded version by truncation.

The endpoint spins are uniform signs.  Since \(f_u,f_v\) are
centered, their conditional expectations have the form
\[
 \E_{\mathrm{pl}}[f_u\mid\sigma_a]=\alpha_u\sigma_a,\qquad
 \E_{\mathrm{pl}}[f_v\mid\sigma_c]=\alpha_v\sigma_c.
\]
\(L^2\)-contraction of conditional expectation gives
\(|\alpha_u|,|\alpha_v|\le\|f\|\).  The broadcast correlation along
the connecting path is
\(\E_{\mathrm{pl}}[\sigma_a\sigma_c]=\theta^{k-2R}\).
It follows that
\[
 |\E_\nu[f_uf_v]|
   =|\alpha_u\alpha_v|\theta^{k-2R}
   \le\|f\|^2 b^R b^{-k/2}.
\]
For \(k\le2R\) the same bound follows from Cauchy--Schwarz, since
\(b^R b^{-k/2}\ge1\).  We have therefore proved
\begin{equation}
\label{eq:local-correlation}
 |\E_\nu[f_uf_v]|
 \le C_f b^{-\dist(u,v)/2}
 \quad\text{for all \(u,v\)},\qquad C_f=b^R\|f\|^2.
\end{equation}
This calculation uses the planted spins to describe the law, while
the assignment itself remains a function of the observations alone.

We now turn spatial decay into an operator bound.  Let \(p_n(o,v)\)
be the \(n\)-step transition probabilities of simple random walk on
the tree and put \(q(v)=b^{-\dist(o,v)/2}\).  For the deterministic
neighbor average \(P\) on functions of vertices,
\[
 (Pq)(v)=\rho q(v)\quad(v\ne o),\qquad
 (Pq)(o)=b^{-1/2}\le\rho q(o).
\]
The first identity uses one neighbor closer to \(o\) and \(b\)
farther away.  Positivity of \(P\) gives \(P^nq\le\rho^nq\) by
induction, and hence
\[
 \sum_vp_n(o,v)q(v)\le\rho^n.
\]
Only finite-step random-walk averages are used here; \(q\) need
not belong to \(\ell^2(\mathcal V)\).
Self-adjointness of \(\mathsf A_\nu\), its finite-walk expansion, and
\eqref{eq:local-correlation} now give, for every \(n\ge0\),
\begin{equation}
\label{eq:marginal-even-moments}
 \|\mathsf A_\nu^nf\|^2
 =\langle f,\mathsf A_\nu^{2n}f\rangle
 =\sum_vp_{2n}(o,v)\E_\nu[f_of_v]
 \le C_f\rho^{2n}.
\end{equation}

Let \(\mu_f\) be the spectral measure of \(f\) for the self-adjoint
restriction of \(\mathsf A_\nu\) to centered assignments.  For
\(\varepsilon>0\), \eqref{eq:marginal-even-moments} implies
\[
 \mu_f\{|\lambda|\ge\rho+\varepsilon\}
 \le C_f\left(\frac{\rho}{\rho+\varepsilon}\right)^{2n}.
\]
Letting \(n\to\infty\) proves that \(\mu_f\) is supported on
\([-\rho,\rho]\).

It remains to pass from local assignments to all centered ones.
For \(f\in\mathcal H_{\mathrm v}(\nu)\) centered, let
\[
 f_v^{(R)}
 =\E_\nu\!\left[f_v\ \middle|\ 
     \sigma\{X_z(s):z\in B_R(v),\ s\ge0\}\right].
\]
These are centered square-integrable local assignments.
To verify their equivariance, pull the conditional expectation at
\(\gamma v\) back by a fixed automorphism \(\gamma\).  Invariance of
\(\nu\) pulls its conditioning sigma-field back to that for \(B_R(v)\),
and equivariance pulls its integrand back to \(f_v\).  Uniqueness of
conditional expectation gives the required identity almost surely
for that \(\gamma,v\).  As \(R\to\infty\), the conditioning
sigma-fields generate the complete observation array, so
\(f_o^{(R)}\to f_o\) in \(L^2\).  This proves density of the local
assignments just used.  Equally, truncating and recentering \(f\)
first gives a bounded local dense family.
The range of the corresponding spectral projection
\(\mathbf 1_{[-\rho,\rho]}(\mathsf A_\nu)\) on the centered space is
closed and contains
every such local assignment.  It is consequently the entire centered
space, proving the lemma.  Finally,
\(\rho<1\) follows from \(d=b+1\) and
\((b+1)^2-4b=(b-1)^2>0\).
\end{proof}

\subsection{Passing the bound to marginal differences}

Fix the joining \(\pi\) in Proposition~\ref{prop:path-variance}.
For a marginal vertex or edge assignment define
\[
 (\delta f)(X,X')=f(X)-f(X').
\]
This is well defined on the marginal \(L^2\) equivalence classes:
each coordinate under \(\pi\) has law \(\nu\).  The same observation
and the triangle inequality give
\(\|\delta f\|_\pi\le2\|f\|_\nu\).  In the joined assignment spaces
put
\[
 \mathcal K=
 \overline{\delta\mathcal H_{\mathrm v}(\nu)}
 \subseteq\mathcal H_{\mathrm v}(\pi),\qquad
 \mathcal D=
 \overline{\delta\mathcal H_{\mathrm e}(\nu)}
 \subseteq\mathcal H_{\mathrm e}(\pi),
\]
where the closures use the respective \(L^2(\pi)\) norms.
These are closed linear spaces.

The vertex neighbor averages for \(\nu\) and \(\pi\) commute with
\(\delta\), since they use the same finite sums over neighboring
vertex indices.  Thus \(\mathsf A_\pi\mathcal K\subseteq\mathcal K\).
The joined \(\mathsf A_\pi\) is self-adjoint, so the same
orthogonality argument as above shows that \(\mathcal K\) is
reducing and its restriction
\(\mathsf A_{\mathcal K}:=\mathsf A_\pi|_{\mathcal K}\) is
self-adjoint.

Subtracting the constant \(\E_\nu f_o\) from a marginal vertex
assignment does not change its difference.  Lemma~\ref{lem:marginal-spectrum}
therefore gives
\begin{equation}
\label{eq:difference-powers}
 \|\mathsf A_\pi^n\delta f\|_\pi
 =\|\delta(\mathsf A_\nu^n(f-\E_\nu f_o))\|_\pi
 \le2\rho^n\|f-\E_\nu f_o\|_\nu
 \qquad(n\ge0).
\end{equation}
Applying the spectral-measure argument from
\eqref{eq:marginal-even-moments} to the joined self-adjoint
operator \(\mathsf A_\pi\) shows that the spectral measure of \(\delta f\) is
supported on \([-\rho,\rho]\).  The constant on the right of
\eqref{eq:difference-powers} may depend on the chosen preimage
\(f\); only its exponential rate matters as \(n\to\infty\).
In particular, no lower comparison between
\(\|\delta f\|_\pi\) and \(\|f-\E_\nu f_o\|_\nu\) is used.
The closed range of the joined spectral projection
\(\mathbf 1_{[-\rho,\rho]}(\mathsf A_\pi)\) contains the closure of
all these differences.  Hence
\begin{equation}
\label{eq:difference-spectrum}
 \|\mathsf A_{\mathcal K}\|\le\rho.
\end{equation}
This is the vertex estimate needed below.  Its domain is
\(\mathcal K\), the closed space generated by marginal differences.

\subsection{From vertices to nonbacktracking edge averages}

A vertex assignment can be read at either endpoint of a directed
edge.  We compare these two lifts and then isolate the part of an
edge assignment orthogonal to both.  On the joined assignment spaces define
the tail and head lifts
\(\mathsf F,\mathsf H:\mathcal H_{\mathrm v}(\pi)\to
\mathcal H_{\mathrm e}(\pi)\), the reversal \(\mathsf J\), and
the forward nonbacktracking average \(\mathsf N\) by
\begin{align*}
 (\mathsf Ff)_{(i,j)}&=f_i,&
 (\mathsf Hf)_{(i,j)}&=f_j,&
 (\mathsf Jg)_{(i,j)}&=g_{(j,i)},\\
 (\mathsf Ng)_{(i,j)}
   &=\frac1b\sum_{\substack{k\sim j\\k\ne i}}g_{(j,k)}.
\end{align*}
Invariance makes the two lifts isometries and makes \(\mathsf J\)
a self-adjoint isometry with \(\mathsf J^2=I\).  Jensen's inequality
makes \(\mathsf N\) a contraction.  The adjoints of the lifts are
\begin{equation}
\label{eq:lift-adjoints}
 (\mathsf F^*g)_v=\frac1d\sum_{u\sim v}g_{(v,u)},\qquad
 (\mathsf H^*g)_v=\frac1d\sum_{u\sim v}g_{(u,v)}.
\end{equation}
Indeed, in the edge inner product we may average the fixed outgoing
edge over all neighbors of its tail, because the tail stabilizer is
transitive on those neighbors.  This proves the first formula.
Interchanging the endpoints first proves the second.  In particular,
\[
 \mathsf F^*\mathsf H=\mathsf H^*\mathsf F=\mathsf A_\pi,\qquad
 \mathsf J\mathsf F=\mathsf H,\qquad
 \mathsf J\mathsf H=\mathsf F.
\]
Direct substitution in the definitions gives
\begin{align}
 \mathsf N&=\frac db\,\mathsf H\mathsf F^*-\frac1b\,\mathsf J,
 &\mathsf N^*&=\frac db\,\mathsf F\mathsf H^*-\frac1b\,\mathsf J,
 \label{eq:edge-operator}\\
 \mathsf N\mathsf F&=\mathsf H,
 &\mathsf N\mathsf H&=\frac db\,\mathsf H\mathsf A_\pi-\frac1b\,\mathsf F.
 \label{eq:edge-lifts}
\end{align}
Thus
\((\mathsf N^*g)_{(i,j)}
=b^{-1}\sum_{k\sim i,\ k\ne j}g_{(k,i)}\)
and \(\mathsf N^*=\mathsf J\mathsf N\mathsf J\).

The same formulas act on the marginal assignment spaces.  They use
finite sums over neighboring indices, selections of endpoint values,
and edge reversal, so they commute with taking marginal differences.
Apply this first on the ranges of \(\delta\), and then
use boundedness to pass to their closures.  We obtain
\[
 \mathsf F\mathcal K,\mathsf H\mathcal K\subseteq\mathcal D,\qquad
 \mathsf F^*\mathcal D,\mathsf H^*\mathcal D\subseteq\mathcal K,
 \qquad \mathsf J\mathcal D=\mathcal D.
\]
Equation~\eqref{eq:edge-operator} now shows that both \(\mathsf N\)
and \(\mathsf N^*\) preserve \(\mathcal D\).  Thus \(\mathcal D\)
is a reducing subspace for the nonbacktracking operator.

Define
\[
 \mathsf R:\mathcal K\oplus\mathcal K\longrightarrow\mathcal D,
 \qquad \mathsf R(f,g)=\mathsf Ff+\mathsf Hg.
\]
The adjoint identities and \eqref{eq:difference-spectrum} give
\begin{align}
 \|\mathsf R(f,g)\|^2
   &=\|f\|^2+\|g\|^2+2\langle f,\mathsf A_{\mathcal K}g\rangle,
   \notag\\
 (1-\rho)(\|f\|^2+\|g\|^2)
   &\le\|\mathsf R(f,g)\|^2
   \le(1+\rho)(\|f\|^2+\|g\|^2).
 \label{eq:lift-gram}
\end{align}
Here \(2|\langle f,\mathsf A_{\mathcal K}g\rangle|
\le\rho(\|f\|^2+\|g\|^2)\).
Because \(\rho<1\), \(\mathsf R\) is injective with closed range
\[
 \mathcal L=\mathsf F\mathcal K+\mathsf H\mathcal K.
\]
In fact, the lower bound makes the preimages of a convergent sequence
in the range Cauchy.  It also gives
\(\|\mathsf R^{-1}\|\le(1-\rho)^{-1/2}\) on \(\mathcal L\),
while the upper bound gives \(\|\mathsf R\|\le(1+\rho)^{1/2}\).

For \(y\in\mathcal D\ominus\mathcal L\), the vectors
\(\mathsf F^*y,\mathsf H^*y\) lie in \(\mathcal K\) and are
orthogonal to every vector in \(\mathcal K\); for example
\(\langle \mathsf F^*y,f\rangle=\langle y,\mathsf Ff\rangle=0\).
Both therefore vanish.  Reversal preserves \(\mathcal D\) and
interchanges the two lifts in \(\mathcal L\), so it preserves
\(\mathcal D\ominus\mathcal L\) as well.  On this complement,
\eqref{eq:edge-operator} reduces to
\begin{equation}
\label{eq:edge-complement}
 \mathsf Ny=-b^{-1}\mathsf Jy,\qquad
 \|\mathsf N^ny\|=b^{-n}\|y\|\quad(n\ge0).
\end{equation}
The space \(\mathcal L\) is also invariant under \(\mathsf N\)
by \eqref{eq:edge-lifts}.  We have split \(\mathcal D\) into
two orthogonal invariant spaces.  The complement already decays
geometrically; it remains to bound powers on the head and tail range.

On \(\mathcal K\oplus\mathcal K\), \eqref{eq:edge-lifts} reads
\[
 \mathsf R^{-1}\mathsf N\mathsf R
 =\mathsf T=
 \begin{pmatrix}
   0&-b^{-1}I\\
   I&(d/b)\mathsf A_{\mathcal K}
 \end{pmatrix}.
\]
This matrix need not be normal.  We bound its powers by its
second-order recurrence.  Put
\[
 \mathsf Z=\sqrt b\,\mathsf T,\qquad
 \mathsf B=\frac d{2\sqrt b}
 \begin{pmatrix}
   \mathsf A_{\mathcal K}&0\\
   0&\mathsf A_{\mathcal K}
 \end{pmatrix}.
\]
The operator \(\mathsf B\) is self-adjoint and
\(\|\mathsf B\|\le1\).  Matrix multiplication gives
\begin{equation}
\label{eq:chebyshev-identity}
 \mathsf Z\mathsf B=\mathsf B\mathsf Z,\qquad
 \mathsf Z^2-2\mathsf B\mathsf Z+I=0.
\end{equation}
Let \(U_{-1}=0\), \(U_0=1\), and
\(U_{m+1}(x)=2xU_m(x)-U_{m-1}(x)\) be the Chebyshev
polynomials of the second kind.  Induction using
\eqref{eq:chebyshev-identity} gives
\[
 \mathsf Z^n
   =U_{n-1}(\mathsf B)\mathsf Z-U_{n-2}(\mathsf B)
   \qquad(n\ge1).
\]
For \(m\ge0\), the trigonometric identity
\[
 U_m(\cos s)=\sum_{j=0}^m e^{\,\mathrm i(m-2j)s}
\]
gives \(\sup_{[-1,1]}|U_m|\le m+1\), including the endpoints.
The spectral theorem applied to the self-adjoint contraction
\(\mathsf B\) therefore yields
\[
 \|\mathsf Z^n\|\le n\|\mathsf Z\|+(n-1)\qquad(n\ge1).
\]
For completeness, the off-diagonal part of \(\mathsf Z\) has norm
\(\sqrt b\), while its only nonzero diagonal block is
\((d/\sqrt b)\mathsf A_{\mathcal K}\), of norm at most \(2\).
Thus \(\|\mathsf Z\|\le\sqrt b+2\).  At the allowed scalar
endpoints \(\pm1\) for \(\mathsf B\), the
characteristic polynomial in \eqref{eq:chebyshev-identity} has a
double root.  The linear factor in \(n\) accounts for this endpoint
behavior.

Using \(\mathsf T^n=b^{-n/2}\mathsf Z^n\) and the two bounds for
\(\mathsf R\) and \(\mathsf R^{-1}\), we conclude that
\(\|\mathsf N^n|_{\mathcal L}\|
\le C_d(n+1)b^{-n/2}\).
Together with \eqref{eq:edge-complement} and the orthogonal invariant
decomposition, this gives
\begin{equation}
\label{eq:edge-power-bound}
 \|\mathsf N^n|_{\mathcal D}\|
 \le C_d(n+1)b^{-n/2}\qquad(n\ge0).
\end{equation}
The same estimate holds for \((\mathsf N^*)^n\) on \(\mathcal D\),
by reversal.  The spectral theorem was used only for the
self-adjoint vertex averages and \(\mathsf B\); the edge
estimate follows from the displayed recurrence.

\subsection{Summing the edge correlations}

Let \(D=\delta g\in\mathcal D\) be the difference in the proposition.
Its mean at every directed edge is zero, since the two marginals are
equal, and \(\|D\|^2=S\).  Fix a directed simple path and write
\(e_i=(v_{i-1},v_i)\).  If \(i<j\), invariance and transitivity
on directed simple paths give
\begin{equation}
\label{eq:edge-lag-correlation}
 \E_\pi[D_{e_i}D_{e_j}]
   =\langle D,\mathsf N^{j-i}D\rangle_{\mathrm e,\pi}.
\end{equation}
Indeed, \(\mathsf N^{j-i}\) averages over the \(b^{j-i}\)
nonbacktracking continuations of a fixed directed edge.  On a tree
each such continuation is simple.  An automorphism carries it to
the corresponding fixed subpath, so every term has the expectation
on the left.  This is a finite average of identities for fixed
automorphisms.

By \eqref{eq:edge-power-bound}, the absolute value of the lag-\(n\)
correlation is at most \(C_d(n+1)b^{-n/2}S\).  Expanding the square
of the sum and grouping equal lags therefore gives
\begin{align*}
 \E_\pi\left[\left(\sum_{i=1}^rD_{e_i}\right)^2\right]
 &=rS+2\sum_{n=1}^{r-1}(r-n)
                 \langle D,\mathsf N^nD\rangle\\
 &\le rS\left(1+2C_d\sum_{n=1}^\infty(n+1)b^{-n/2}\right)
 \le C'_d rS.
\end{align*}
The series converges because \(b\ge2\).  If all edges are reversed,
order them as
\((v_r,v_{r-1}),\ldots,(v_1,v_0)\).  They are the forward edges
of the reversed simple path, and their sum is unchanged by this
reordering.  The same argument proves the second inequality of
\eqref{eq:path-variance}.

Finally, take \(\pi\) to be the \((X,X')\)-marginal of the joining
in Proposition~\ref{prop:compatible-copies}, forgetting its common
innovation coordinate.  At a fixed finite time \(t\), use
\[
 g_{(i,j)}(X)=\log a(h_{j\to i}(t)).
\]
The component versions constructed earlier make this an equivariant
marginal assignment.  The cavity field bound and
\((\log a)(0)=0\), \(|(\log a)'|\le2\), imply
\[
 |g_{(i,j)}(X)|
 \le2|h_{j\to i}(t)|
 \le2\bigl(|X_j(t)|+b\beta\bigr).
\]
Since \(\E_\nu X_j(t)^2=t+t^2\), it is square integrable.
Its difference is precisely \(D_{(i,j)}(t)\).  Applying the general
estimate completes the proof of Proposition~\ref{prop:path-variance}.
The bound is uniform in the time and in the joining because its
constant depends only on \(d\).
